%=============================================================================!
% 4.0  CHAPTER OPENING                                                        !
%=============================================================================!
Chapter~2 solved the motion of a block on a spring, and Chapter~3 showed
that near a smooth potential-energy minimum with positive curvature the
restoring force is approximately that of a spring. This connects an
ordinary spring to the vibrations of atoms near a stable arrangement.
We now examine how the motion depends on the shape of the well and on
the forces supplied by the surroundings.

For the harmonic spring, circular motion, the exchange of kinetic and
potential energy, and the path through the plane of states give three
ways to follow the same oscillation. Larger displacements explore terms
beyond the quadratic approximation, so the period can depend on the
energy; in the pendulum and pair potentials studied here, it increases.
Drag removes energy and damps the motion, while a periodic force can
supply energy most effectively near resonance. These effects return in
thermostats and in the response of molecules to light.

When several atoms move, their displacements are coupled because moving
one changes the forces on its neighbours. We first reduce the two-body
problem to a single relative coordinate, then find the independent
patterns of a coupled system, its normal modes, using the eigenvalues
of a matrix. Molecular examples connect these frequencies to periods of
less than ten femtoseconds, establishing the scale that a simulation's
time step must resolve. Chapter~5 then turns to rotation.
